\documentclass{article}
\usepackage[utf8]{inputenc}
\usepackage{amsmath,amssymb,amsthm,amsfonts}
\usepackage{geometry}
\geometry{letterpaper, margin=1in}

\title{Quantifying the Effective Sample Size of Task Descriptors in In-Context Learning}
\author{Bruce Quan Nguyen}
\date{\today}

\newtheorem{definition}{Definition}
\newtheorem{proposition}{Proposition}

\begin{document}

\maketitle

\begin{abstract}
In-context learning (ICL) allows models to adapt to tasks using prompts that typically combine a natural language task descriptor with input-output demonstrations. While the informational value of demonstrations is well understood, quantifying the precise contribution of the task descriptor remains challenging. By formulating ICL as Bayesian inference over a latent task, we treat the descriptor and demonstrations as distinct modalities of observations. We quantify the descriptor's value via its \emph{effective sample size} (ESS)---the equivalent number of demonstrations required to achieve an identical reduction in one-step predictive log-loss regret for a Bayes-optimal predictor. In the exact setting of in-context linear regression, we demonstrate that a descriptor's ESS is governed by its \emph{precision} (the reduction in prior variance) and its \emph{reliability} (the probability of correct specification). We prove that for a correctly specified descriptor, ESS scales linearly with precision and task dimension. Conversely, for a potentially misspecified descriptor, the ESS is strictly bounded without demonstrations, as the optimal predictor must maintain posterior uncertainty over the descriptor's validity. Demonstrations act as evidence to resolve this unidentifiability, effectively unlocking the descriptor's value. Furthermore, we show that small meta-trained Transformers behave Bayes-optimally within the constraints of their output representations; when faced with misspecified descriptors, unimodal output heads force the model to project the optimal bimodal mixture onto a single Gaussian. Finally, we demonstrate that a network's reliance on the descriptor is entirely determined by its meta-training distribution, where over-reliance on highly precise but misspecified descriptors yields catastrophic regret.
\end{abstract}

\section{Introduction}
The empirical success of large language models is heavily reliant on in-context learning (ICL), where a model adapts to a novel task given a prompt. Standard prompts incorporate two distinct sources of task information: a natural language task descriptor and a few demonstration examples. Empirical observations regarding the descriptor's utility are heavily mixed; while a well-crafted descriptor can substitute for hundreds of examples, models are also known to perform adequately with misleading descriptors or heavily prioritize demonstrations over explicit instructions. Existing theoretical frameworks for ICL often model the prompt strictly as a sequence of demonstrations, leaving a fundamental theoretical gap: how can we rigorously quantify the informational equivalent of a task descriptor?

We address this by adopting a Bayesian perspective on ICL. In this framework, the pretraining distribution constitutes a prior over a latent task. The task descriptor acts as a direct, explicit observation of the task parameters, whereas demonstrations serve as indirect observations. Conditioning on the descriptor yields an updated task prior. In Bayesian statistics, the information contained in a prior can be quantified as its effective sample size (ESS), which is dictated by the prior's concentration (its \emph{precision}) and the probability of prior-data conflict (its \emph{reliability}).

We analyze this dynamic within noisy in-context linear regression with a Gaussian task prior, where Bayes-optimal inference admits closed-form solutions. This tractability allows us to exactly compute a descriptor's ESS and provides exact reference predictors to evaluate the behavior of meta-trained Transformer models.

\section{Related Work}
Our approach builds upon the understanding of ICL as implicit Bayesian inference, where meta-trained Transformers have been shown to implement optimal estimators such as Bayes-optimal ridge regression. While some literature has integrated task descriptors---modeling them as tokens that indicate hypothesis classes or carry input means---these works generally do not quantify the descriptor's value in units of examples, nor do they formalize the probability of descriptor misspecification. Our ESS formulation aligns with prior sample size methodologies in Bayesian statistics, particularly those utilizing predictive regret. To model potential misspecification, we employ robust mixture priors, a technique commonly used in clinical trials to integrate potentially conflicting historical data.

\section{Problem Formulation}

\paragraph{Tasks and Demonstrations.} We assume a latent task represented by a weight vector $w \in \mathbb{R}^d$. A demonstration consists of an input $x \sim \mathcal{N}(0, I_d)$ and a label $y = w^\top x + \epsilon$, with observation noise $\epsilon \sim \mathcal{N}(0, \sigma_y^2)$. The base prior over tasks is isotropic Gaussian $w \sim \mathcal{N}(0, s_0^2 I_d)$. The prior signal-to-noise ratio (SNR) is defined as $a_0 = s_0^2/\sigma_y^2$, which we set to $a_0 = 10$.

\paragraph{Task Descriptors.} A descriptor is formalized as a vector $m \in \mathbb{R}^d$. Its generative process depends on a latent Bernoulli indicator $Z \sim \text{Bernoulli}(p)$, which dictates whether the descriptor is correctly specified ($Z=1$) or misspecified ($Z=0$). The generative process is:
\begin{align}
    m &\sim \mathcal{N}(0, (1-r)s_0^2 I_d) \\
    w \mid m, Z=1 &\sim \mathcal{N}(m, r s_0^2 I_d) \\
    w \mid m, Z=0 &\sim \mathcal{N}(0, s_0^2 I_d)
\end{align}
Marginalizing over $Z$, the prior for $w$ given the descriptor $m$ is the robust mixture $p \mathcal{N}(m, rs_0^2 I_d) + (1-p) \mathcal{N}(0, s_0^2 I_d)$. The parameter $r \in (0, 1)$ denotes the descriptor's \textbf{precision} (the ratio of the conditional variance to the base prior variance). The parameter $p$ denotes the descriptor's \textbf{reliability}.

\paragraph{Regret and ESS.} We evaluate a predictor's uncertainty given a prompt $D$ (the descriptor, $n$ examples, and a query $x$) via its one-step expected excess log-loss regret relative to an oracle that knows the true task $w$:
\begin{equation}
    R = \mathbb{E} \left[ -\log q(y \mid D) + \log \phi(y; w^\top x, \sigma_y^2) \right]
\end{equation}
where $\phi$ is the Gaussian density. Let $R^{\text{desc}}(n)$ denote the Bayes-optimal regret with the descriptor and $n$ examples, and $R^{\text{ex}}(n)$ the regret with $n$ examples alone. 
\begin{definition}
The Effective Sample Size (ESS) of a descriptor, given $n$ demonstrations, is $\text{ESS}_n = n^* - n$, where $n^*$ is the root of $R^{\text{ex}}(n^*) = R^{\text{desc}}(n)$, calculated via linear interpolation. The standalone ESS is $\text{ESS}_0$.
\end{definition}

\paragraph{Reference Predictors and Reliance.} A learner applies an effective mixture weight $q$ to the descriptor component. An \emph{oracle-calibrated} predictor uses $q=p$ (Bayes-optimal). An \emph{overconfident} predictor assumes $q=1$. We define a learner's \emph{implied reliance} as the value $q$ that minimizes the mean squared error between its posterior-mean predictions and those of the corresponding reference model.

\section{The Value of Precision (Correctly Specified Descriptors)}
When a descriptor has perfect reliability ($p=1$), the posterior remains Gaussian. The regret is half the expected log of the remaining predictive variance, $R = \frac{1}{2}\mathbb{E} \log(1 + x^\top \Sigma x)$.

\begin{proposition} \label{prop:reliable}
For $p=1$ and conjugate-prior ESS $c = 1/(ra_0) = \sigma_y^2 / (rs_0^2)$:
\begin{enumerate}
    \item (Coarse descriptors) If $c \le 1$, then as $rd \to \infty$, $\text{ESS} = (d-1)(1-r) + (1-r)c + O(1/(rd))$.
    \item (Precise descriptors) If $c \ge d$, then as $c/d \to \infty$, $\text{ESS} = c + d + 1 - 1/a_0 + O(d/c)$.
    \item For every $d, a_0, r$, we have $\text{ESS} \le c + d + 2$.
\end{enumerate}
\end{proposition}
Both asymptotic regimes demonstrate that the ESS of a correctly specified descriptor decomposes into a high-SNR term (representing the fraction of the dimension it constrains) and the conjugate-prior ESS $c$. Roughly, $\text{ESS} \approx (d-1)(1-r) + 1/(ra_0)$.

\section{The Identifiability Bottleneck (Potentially Misspecified Descriptors)}
When $p < 1$, the Bayes-optimal posterior predictive is a two-component mixture. The optimal predictor must maintain uncertainty over the validity of the descriptor.

\begin{proposition} \label{prop:unreliable}
Let $\pi_n(Z)$ denote the posterior weight on the true mixture component. The regret is bounded by:
\begin{equation}
    p R^{\text{desc}}_{p=1}(n) + (1-p) R^{\text{ex}}(n) \le R^{\text{desc}}(n) \le p R^{\text{desc}}_{p=1}(n) + (1-p) R^{\text{ex}}(n) + \mathbb{E}[-\log \pi_n(Z)]
\end{equation}
\end{proposition}
The identification cost $\mathbb{E}[-\log \pi_n(Z)]$ is the conditional entropy $H(Z \mid D_n)$. Without demonstrations ($n=0$), the regret is bounded below by $(1-p)R^{\text{ex}}(0)$. Consequently, the ESS strictly saturates at a finite value, preventing arbitrary scaling with precision as $r \to 0$.

\paragraph{Demonstrations resolve unidentifiability.} As demonstrations accumulate, they provide evidence to infer the latent indicator $Z$, driving the conditional entropy $H(Z \mid D_n) \to 0$. The ESS of a misspecified descriptor then approaches $p$ times that of a correctly specified one ($\text{ESS}_n \to p c$). Demonstrations effectively pay down the informational cost of potential misspecification.

\paragraph{The cost of overconfidence.} An overconfident predictor ($q=1$) interacting with a misspecified descriptor incurs severe regret. Because the assumed prior is highly concentrated around the erroneous descriptor, the true generative target falls drastically outside the posterior predictive distribution. For highly precise descriptors, overconfidence is an active liability, yielding higher regret than ignoring the descriptor entirely. 

\section{Empirical Analysis: Meta-Trained Transformers}
We meta-trained 12-layer Transformers on sequences generated by our problem formulation ($d=5, a_0=10$). The descriptor $m$ is provided via a dedicated prefix token, optimized over Gaussian log-loss.

\paragraph{Expressivity bounds Bayes-optimal behavior.} On their respective training distributions, Transformers match the Bayes-optimal ESS for perfectly specified descriptors. However, for potentially misspecified descriptors with zero demonstrations, the networks fall short of the exact Bayes-optimal mixture ESS. Standard regression heads emit a single mean and variance, inherently lacking the capacity to represent a bimodal posterior mixture. Consequently, the networks project the optimal mixture onto the \emph{best single Gaussian} via moment matching. As demonstrations collapse the posterior mixture to a single component, the networks successfully converge with the exact Bayes-optimal predictor.

\paragraph{Reliance is inherited from meta-training.} We demonstrate that the implied reliance of the network is governed strictly by the meta-training distribution. Models trained exclusively on perfect reliability ($p=1$) learn an implied reliance of $q \approx 1$. At deployment, they act as overconfident predictors, suffering catastrophic regret on precise but misspecified descriptors. Conversely, models trained on $p=0.9$ learn to maintain structural uncertainty, exhibiting an implied reliance of $q \approx 0.85$ even when evaluated exclusively on correctly specified descriptors.

\section{Discussion and Limitations}
Our theoretical framework establishes that the ESS of a valid task descriptor scales with its precision and constrained dimensionality. However, potential misspecification introduces an identifiability bottleneck, forcing optimal predictors to maintain uncertainty and capping the descriptor's value. Demonstrations are critical not just as additive data, but as verification signals that resolve this unidentifiability. Meta-trained Transformers approximate this Bayes-optimal reasoning, bounded by the expressivity of their unimodal output representations, with their propensity for reliance statically inherited from training.

Our analysis is subject to limitations, notably the assumption of a linear-Gaussian generative process. Furthermore, our ESS metric is defined via one-step prediction; under sequential prediction horizons where the predictor's own feedback accumulates over time, the ESS dynamics shift slightly, representing a valuable direction for future research.

\appendix
\section{Mathematical Proofs}

\subsection{Proof of Proposition \ref{prop:reliable} (Correctly Specified Descriptors)}
We prove Proposition \ref{prop:reliable} using the exact regret expansions for the predictive log-loss. 
Let $\psi$ denote the digamma function, and let $\varepsilon = 1/a_0$. Given $n$ examples $X_n \in \mathbb{R}^{n \times d}$, the posterior covariance is $\Sigma = \sigma_y^2 (\varepsilon I + X_n^\top X_n)^{-1}$.
For a query $x \sim \mathcal{N}(0, I_d)$, let $A = \|x\|^2 \sim \chi^2_d$. The one-step expected regret is $R^{\text{ex}}(n) = \frac{1}{2} \mathbb{E} \log(1 + x^\top \Sigma x)$.

\paragraph{Regret Expansions.}
By leveraging properties of Wishart distributions and the Sherman-Morrison formula, we can bound the regret in two regimes:
\begin{enumerate}
    \item \textbf{Fewer than $d$ examples ($n < d$):} Let $k = d - n$. The regret is lower bounded by the trace of the inverse Wishart components:
    \begin{equation}
        2R^{\text{ex}}(n) \ge \log(2a_0) + \psi(k/2)
    \end{equation}
    with equality as $a_0 \to \infty$.
    \item \textbf{At least $d$ examples ($n \ge d$):} Let $m = n - d + 1$. The regret is upper bounded by:
    \begin{equation}
        2R^{\text{ex}}(n) \le \psi\left(\frac{n+1}{2}\right) - \psi\left(\frac{m}{2}\right)
    \end{equation}
\end{enumerate}

\paragraph{Locating the ESS Root.}
By Definition 1, the ESS is the root $n^*$ of $R^{\text{ex}}(n^*) = R^{\text{desc}}(0)$. For a correctly specified descriptor, $R^{\text{desc}}(0) = \frac{1}{2}\mathbb{E}\log(1 + c^{-1} A)$ where $c = 1/(ra_0)$.
Matching the regret expansions to $R^{\text{desc}}(0)$ and applying the asymptotic expansions of the digamma function ($\psi(z) = \log z - 1/(2z) + O(z^{-2})$) yields the roots:
\begin{itemize}
    \item In the coarse regime ($c \le 1$), matching the lower bound gives $d - n \approx rd$, leading to $\text{ESS} \approx (d-1)(1-r) + (1-r)c$.
    \item In the precise regime ($c \ge d$), matching the upper bound gives $n \approx c + d + 1$, leading to $\text{ESS} \approx c + d + 1 - 1/a_0$.
\end{itemize}
By concavity of the log-loss, the ESS is strictly bounded across all regimes by $c + d + 2$, concluding the proof.

\subsection{Proof of Proposition \ref{prop:unreliable} (Potentially Misspecified Descriptors)}
Let $D_n = (x_{1:n}, y_{1:n}, m, x)$ denote the prompt and the query input. Let $Z \in \{0, 1\}$ be the true validity indicator of the descriptor.

\paragraph{Step 1: The Mixture Predictive.}
The Bayes-optimal posterior predictive is a $\pi_n$-mixture of the correctly specified predictive $f_1$ and the base prior predictive $f_0$, where $\pi_n = P(Z=1 \mid D_n)$. 

\paragraph{Step 2: The Oracle Twin Lower Bound.}
Consider an oracle predictor that is explicitly given $Z$. This twin predicts using $f_1$ when $Z=1$ (incurring regret $R^{\text{desc}}_{p=1}(n)$) and $f_0$ when $Z=0$ (incurring regret $R^{\text{ex}}(n)$). The twin's expected regret is exactly $p R^{\text{desc}}_{p=1}(n) + (1-p) R^{\text{ex}}(n)$.
Because additional information cannot hurt a Bayes-optimal predictor, the expected log-loss of our mixture predictor $q$ cannot be lower than the oracle twin's log-loss. Thus:
\begin{equation}
    R^{\text{desc}}(n) \ge p R^{\text{desc}}_{p=1}(n) + (1-p) R^{\text{ex}}(n)
\end{equation}

\paragraph{Step 3: The Identification Cost Upper Bound.}
Pointwise, the mixture density satisfies $q \ge \pi_n f_1$ and $q \ge (1-\pi_n) f_0$. Consequently:
\begin{align}
    -\log q &\le -\log f_1 - \log \pi_n \\
    -\log q &\le -\log f_0 - \log(1-\pi_n)
\end{align}
Taking the expectation over the true data generating process $f_Z$ and the latent variable $Z$ yields the upper bound. The penalty term simplifies to the conditional entropy $\mathbb{E}[-\log \pi_n(Z)] = H(Z \mid D_n)$.

\paragraph{Step 4: Monotonicity and Saturation.}
The conditional entropy $H(Z \mid D_n)$ is non-increasing in $n$, since conditioning on additional demonstrations can only reduce entropy. At $n=0$, $H(Z \mid D_0) = H(p)$. 
Crucially, evaluating the lower bound at $n=0$ gives $R^{\text{desc}}(0) \ge (1-p)R^{\text{ex}}(0)$. Since $R^{\text{ex}}(n)$ is strictly decreasing in $n$, there exists a finite $n^*$ such that $R^{\text{ex}}(n^*) = (1-p)R^{\text{ex}}(0)$. The ESS cannot exceed this $n^*$, forcing the strict saturation bottleneck.

\end{document}
